\documentclass[11pt]{article}

\usepackage{amsmath,amssymb,amsthm}
\usepackage{geometry}
\geometry{margin=1in}

\newtheorem{theorem}{Theorem}
\newtheorem{proposition}[theorem]{Proposition}
\newtheorem{lemma}[theorem]{Lemma}
\newtheorem{corollary}[theorem]{Corollary}
\theoremstyle{definition}
\newtheorem{definition}[theorem]{Definition}
\newtheorem{remark}[theorem]{Remark}
\newtheorem{example}[theorem]{Example}

\DeclareMathOperator{\Sym}{Sym}
\DeclareMathOperator{\Stab}{Stab}
\DeclareMathOperator{\Orb}{Orb}
\DeclareMathOperator{\vol}{vol}
\DeclareMathOperator{\sdf}{sdf}

\title{$G$-Carving Grammars:\\
Monotone Descent and a Finite Classification of\\
Closed Paths in the Lattice of $G$-Invariant SDFs}
\author{}
\date{}

\begin{document}
\maketitle

\begin{abstract}
Let $G$ be a finite subgroup of $O(3)$ and let $\mathcal{L}_G$ denote the
lattice of $G$-invariant signed distance functions on $\mathbb{R}^3$ under
pointwise $\min$ and $\max$. We study finite sequences of carving operations
$(\text{subtract}, \text{intersect})$ applied to a canonical $G$-invariant
seed, with each operation using a symmetrised child primitive placed in the
fundamental domain $\mathrm{FD}(G)$. We prove four small structural results:
(i) every carving grammar defines a monotone descending chain of solids;
(ii) the orbit-size function on $\mathrm{FD}(G)$ is upper semicontinuous and
stratifies the domain by stabilizer rank; (iii) the space of $n$-step
``resonant'' grammars, obtained by restricting $(\varphi, r)$ to
stabilised strata and canonical radii, is finite up to continuous shape
parameters; (iv) every such chain converges in the Hausdorff metric to a
$G$-invariant closed limit. These statements formalise the notion that a
carving grammar is a closed path from seed to result inside $\mathcal{L}_G$,
and they give a combinatorial classification of resonant grammars that
matches the discrete vocabulary used in the companion evolver.
\end{abstract}

\section{Setup and notation}
\label{sec:setup}

Let $G \subset O(3)$ be finite. For this paper we have in mind
$G \in \{T_d, O_h, I_h\}$ of orders $24$, $48$, $120$, corresponding
respectively to the tetrahedral, octahedral, and icosahedral symmetry
groups, but the statements hold for any finite subgroup of $O(3)$
acting linearly on $\mathbb{R}^3$.

\begin{definition}[$G$-invariant SDF]
A function $f : \mathbb{R}^3 \to \mathbb{R}$ is a \emph{$G$-invariant signed
distance function} if
\begin{enumerate}
\item $|f(x) - f(y)| \le \|x-y\|$ (1-Lipschitz),
\item $f(g\cdot x) = f(x)$ for every $g \in G$ and $x \in \mathbb{R}^3$.
\end{enumerate}
The collection of such functions is denoted $\mathcal{L}_G$. The
associated solid is $S_f := \{x \in \mathbb{R}^3 : f(x) \le 0\}$, which
is $G$-invariant as a set.
\end{definition}

\begin{lemma}[Lattice structure]
\label{lem:lattice}
The pair $(\mathcal{L}_G, \le)$, with pointwise ordering $f \le g \iff
f(x) \le g(x)\; \forall x$, is a complete lattice under
\[
  (f \wedge g)(x) := \min(f(x), g(x)), \qquad
  (f \vee g)(x) := \max(f(x), g(x)).
\]
Moreover $\min$ and $\max$ preserve both properties $(1)$ and $(2)$, so
$\mathcal{L}_G$ is closed under them.
\end{lemma}

\begin{proof}
1-Lipschitz is preserved because $\min$ and $\max$ of 1-Lipschitz functions
are 1-Lipschitz. $G$-invariance is preserved since for any $g \in G$,
$\min(f(gx), f'(gx)) = \min(f(x), f'(x))$. Completeness follows from the
lattice of bounded 1-Lipschitz functions being complete under pointwise
$\sup$ and $\inf$, and the $G$-invariant subspace being closed under
arbitrary pointwise meets and joins.
\end{proof}

Write $S_G \subset \mathbb{R}^3$ for a compact $G$-invariant \emph{seed}
solid, with canonical SDF $f_0 \in \mathcal{L}_G$ (e.g.\ the inscribed
sphere, tetrahedron, cube, octahedron, dodecahedron, or icosahedron,
each canonically positioned at the origin so that $G$ acts linearly).

\subsection{Fundamental domain and orbit symmetrisation}

The action of $G$ on the unit sphere $S^2$ admits a spherical
\emph{fundamental domain} $\mathrm{FD}(G) \subset S^2$: a closed
spherical triangle with area $4\pi/|G|$ whose $G$-translates tile $S^2$
with overlap only on a measure-zero set (the edges of $\mathrm{FD}$,
which lie on rotation axes). For $T_d$, $O_h$, $I_h$, $\mathrm{FD}(G)$
is the spherical triangle with vertices at a $3$-fold axis (vertex of
the dual polyhedron), a $2$-fold axis (edge midpoint), and a $k$-fold
axis ($k=3,4,5$, face centre), respectively.

\begin{definition}[Symmetrisation]
For $\varphi \in \mathbb{R}^3$ and a primitive SDF $P : \mathbb{R}^3 \to
\mathbb{R}$, define
\[
  \Sym_G(P, \varphi)(x) \;:=\; \min_{g \in G}\, P(g\cdot x - \varphi).
\]
\end{definition}

\begin{proposition}[$G$-invariance of symmetrisation]
\label{prop:ginv}
For any primitive $P$ and any placement $\varphi$,
$\Sym_G(P, \varphi) \in \mathcal{L}_G$ (provided $P$ is 1-Lipschitz).
\end{proposition}

\begin{proof}
For $h \in G$,
$\min_{g\in G} P(g\cdot hx - \varphi) = \min_{g' \in G} P(g' x - \varphi)$
by the substitution $g' = gh$, using that $G$ is closed under right
multiplication. 1-Lipschitz descends from $P$ through
$x \mapsto g\cdot x$ (isometry) and pointwise minimum.
\end{proof}

\subsection{Smooth Booleans as lattice operations}

Define, for $k \ge 0$,
\[
  \mathrm{smax}_k(a,b) := \tfrac12(a+b) + \tfrac12|a - b| + \mathrm{bump}_k(a,b),
\]
where $\mathrm{bump}_k \ge 0$ vanishes for $k=0$ and satisfies
$\mathrm{bump}_k \to 0$ pointwise as $k \to 0^+$. The smooth operation
$\mathrm{smax}_k$ is an upper bound on $\max$: $\mathrm{smax}_k(a,b) \ge
\max(a,b)$ for all $a,b$, with equality at $k=0$. Analogously
$\mathrm{smin}_k(a,b) \le \min(a,b)$.

\begin{lemma}
\label{lem:smax-monotone}
$\mathrm{smax}_k(f, g)(x) \ge f(x)$ pointwise. Hence if
$h := \mathrm{smax}_k(f, g)$, then $S_h = \{h \le 0\} \subseteq S_f$.
Symmetrically $\mathrm{smin}_k(f, g) \le f$ and
$S_{\mathrm{smin}_k(f,g)} \supseteq S_f$.
\end{lemma}

\begin{proof}
Immediate from $\mathrm{smax}_k(a,b) \ge \max(a,b) \ge a$, and the
pointwise containment of sublevel sets reverses pointwise order.
\end{proof}

\section{$G$-carving grammars}
\label{sec:grammars}

\begin{definition}[$G$-carving grammar]
A \emph{$G$-carving grammar of length $n$} is a tuple
\[
  \Gamma = (f_0;\; (b_1, P_1, \varphi_1, r_1, k_1), \dots,
           (b_n, P_n, \varphi_n, r_n, k_n)),
\]
where
\begin{itemize}
\item $f_0 \in \mathcal{L}_G$ is a canonical seed;
\item $b_i \in \{\mathsf{subtract}, \mathsf{intersect}\}$;
\item $P_i$ is a 1-Lipschitz primitive SDF (sphere, tetrahedron, cube,
      icosahedron);
\item $\varphi_i \in \mathrm{FD}(G) \subset S^2$ and $r_i > 0$ specify
      direction and distance of the child centre;
\item $k_i \ge 0$ is a smoothing radius.
\end{itemize}
The associated sequence $(f_i)_{i=0}^{n}$ is defined by
\[
  c_i(x) \;:=\; \Sym_G(P_i, r_i \varphi_i)(x),
\]
\[
  f_i(x) \;:=\;
  \begin{cases}
    \mathrm{smax}_{k_i}(f_{i-1}(x),\, -c_i(x)) & \text{if } b_i = \mathsf{subtract},\\
    \mathrm{smax}_{k_i}(f_{i-1}(x),\, c_i(x)) & \text{if } b_i = \mathsf{intersect}.
  \end{cases}
\]
The associated sequence of solids is $S_i := \{f_i \le 0\}$.
\end{definition}

\begin{remark}
Omitting $\mathsf{add}$ (smooth union) is the defining feature of a
\emph{carving} grammar. The union operator would use
$\mathrm{smin}_k(f_{i-1}, c_i)$ and therefore grow the solid --- it is
structurally incompatible with the monotone descent property proved
below.
\end{remark}

\section{Monotone descent}
\label{sec:descent}

\begin{theorem}[Monotone descent]
\label{thm:monotone}
Let $\Gamma$ be any $G$-carving grammar of length $n$. Then
\[
  S_0 \;\supseteq\; S_1 \;\supseteq\; S_2 \;\supseteq\; \cdots
  \;\supseteq\; S_n,
\]
and each $S_i$ is $G$-invariant as a set.
\end{theorem}

\begin{proof}
By Proposition~\ref{prop:ginv}, $c_i \in \mathcal{L}_G$.
By Lemma~\ref{lem:lattice}, $\mathcal{L}_G$ is closed under pointwise
$\max$; the smooth version $\mathrm{smax}_{k_i}$ is an over-approximation
that still lies in $\mathcal{L}_G$ (it commutes with $G$-action because
$f_{i-1}$ and $c_i$ are $G$-invariant). So $f_i \in \mathcal{L}_G$ for
every $i$, and $S_i$ is $G$-invariant.

For the inclusion, Lemma~\ref{lem:smax-monotone} gives $f_i \ge f_{i-1}$
pointwise (whether $b_i$ is subtract or intersect, both branches feed
$f_{i-1}$ as the first argument of $\mathrm{smax}_{k_i}$). Taking
sublevel sets reverses the inequality: $S_i \subseteq S_{i-1}$.
\end{proof}

\begin{corollary}[Closed path interpretation]
Every carving grammar defines a monotone path
\[
  f_0 \;\le\; f_1 \;\le\; \cdots \;\le\; f_n
\]
in the poset $(\mathcal{L}_G, \le)$, starting at the seed and ending at
the result. The path stays inside $\mathcal{L}_G$ throughout, and the
associated solid path $S_0 \supseteq \cdots \supseteq S_n$ is a
descending chain of $G$-invariant compact sets.
\end{corollary}

\section{Stratification of $\mathrm{FD}(G)$ by orbit size}
\label{sec:strat}

\begin{definition}[Orbit size function]
For $\varphi \in S^2$, write $\Stab_G(\varphi) := \{g \in G : g\varphi
= \varphi\}$ and $\Orb_G(\varphi) := \{g\varphi : g \in G\}$. By the
orbit--stabiliser theorem,
\[
  |\Orb_G(\varphi)| \;=\; |G| / |\Stab_G(\varphi)|.
\]
Let $N_G(\varphi) := |\Orb_G(\varphi)|$.
\end{definition}

\begin{theorem}[Stabiliser stratification]
\label{thm:strat}
The function $N_G : S^2 \to \mathbb{N}$ is upper semicontinuous. It
takes the maximum value $|G|$ on a dense open set, and drops on a
finite union of great-circle arcs (fixed sets of order-$2$ rotations)
and a finite set of isolated points (fixed by order $\ge 3$ rotations).
On $\mathrm{FD}(G)$ specifically, the stratification reads:
\begin{itemize}
\item Interior of $\mathrm{FD}(G)$: $N_G = |G|$.
\item Relative interior of edges of $\mathrm{FD}(G)$: $N_G = |G|/2$.
\item Vertices of $\mathrm{FD}(G)$: $N_G = |G|/p$ with
      $p \in \{3, 4, 5\}$ depending on the axis order at that vertex.
\end{itemize}
\end{theorem}

\begin{proof}[Proof sketch]
Upper semicontinuity follows from the fact that
$\{\varphi : |\Stab_G(\varphi)| \ge s\}$ is closed for every $s$: it is
the union of finitely many fixed sets of subgroups $H \le G$, each of
which is a union of great circles (for $H$ generated by a single
reflection or rotation axis) or isolated points (for $H$ containing
multiple independent rotations). On $\mathrm{FD}(G)$, the fundamental
spherical triangle meets each rotation-axis locus exactly on its
boundary: edges correspond to $C_2$ axes (stabiliser order 2), and the
three vertices correspond to the three distinguished rotation axes of
$G$, of orders $3$, $2$, and $p$ (the highest-order axis of $G$).
\end{proof}

\begin{example}[Octahedral $O_h$, $|G|=48$]
Generic $\varphi$: orbit size $48$; edge midpoints of $\mathrm{FD}$
(lying on $C_2$ axes): orbit size $24$; vertices of $\mathrm{FD}$: $8$
($C_3$ vertex axis), $12$ ($C_2$ edge axis), $6$ ($C_4$ face axis).
\end{example}

\begin{corollary}
\label{cor:reduced-orbit}
If $\varphi$ lies on a stratum with $|\Stab_G(\varphi)| = s$, then
$\Sym_G(P, \varphi)(x) = \min_{g \in G/\Stab_G(\varphi)} P(g\cdot x -
\varphi)$, a minimum over $|G|/s$ terms rather than $|G|$. The
symmetrised child is correspondingly supported by fewer orbit copies.
\end{corollary}

\begin{remark}
Placing carves at stabilised strata (edges and vertices of
$\mathrm{FD}(G)$) produces visually and structurally simpler shapes ---
fewer, larger removed features instead of many small ones. This is
exactly the difference between the ``continuous'' and ``resonant''
modes of the companion evolver.
\end{remark}

\section{Finite classification of resonant grammars}
\label{sec:resonant}

\begin{definition}[Resonant grammar]
A carving grammar is \emph{resonant} if for every step $i$,
\begin{itemize}
\item $\varphi_i$ lies in the finite set
    $C_G := V(\mathrm{FD}) \cup E_{\text{mid}}(\mathrm{FD})$
    of vertices and edge midpoints of $\mathrm{FD}(G)$;
\item $r_i \in R_G := \{R, \rho, \iota\}$, the circumradius,
    midradius, and inradius of the seed solid $S_G$;
\item $b_i \in B := \{\mathsf{subtract}, \mathsf{intersect}\}$;
\item $P_i \in \mathcal{P}$, a fixed finite palette of primitive SDFs.
\end{itemize}
\end{definition}

\begin{theorem}[Finite classification]
\label{thm:finite}
Fix $n, G, S_G, \mathcal{P}$. The number of combinatorial types of
$n$-step resonant grammars on $S_G$, where two grammars are equivalent
iff they agree in $(b_i, P_i, \varphi_i, r_i)$ for every $i$, is at
most
\[
  N_{\mathrm{res}}(n) \;\le\;
  \bigl(|B| \cdot |\mathcal{P}| \cdot |C_G| \cdot |R_G|\bigr)^{n}.
\]
In particular, $N_{\mathrm{res}}(n) < \infty$ for every $n$. The
continuous shape parameters (scale factor, smoothing radius $k_i$)
define a bounded open subset of $\mathbb{R}^{2n}$ fibred over this
finite set.
\end{theorem}

\begin{proof}
Each step is a point in the Cartesian product
$B \times \mathcal{P} \times C_G \times R_G$, which is finite. A
length-$n$ grammar is a word of length $n$ in this alphabet, giving at
most $(|B||\mathcal{P}||C_G||R_G|)^n$ distinct words. Continuous
parameters $(s_i, k_i) \in [0,1]^2$ refine each word to a bounded open
region.
\end{proof}

\begin{example}
For the evolver's default vocabulary
($|B|=2$, $|\mathcal{P}|=4$, $|C_G|=6$, $|R_G|=3$), and $n \le 5$, we
have
\[
  N_{\mathrm{res}}(\le 5) \;=\;
  \sum_{i=1}^{5}(2\cdot 4 \cdot 6 \cdot 3)^{i} \;=\;
  \sum_{i=1}^{5} 144^{i} \;<\; 6.2 \times 10^{10}.
\]
This is finite but large. However the fitness filter (monotone descent
plus the carving hard gates: watertight, single-component, minimum wall
thickness) eliminates the vast majority, and --- modulo $G$-action and
continuous rescaling --- the count of \emph{geometrically} distinct
resonant shapes is conjecturally much smaller.
\end{example}

\section{Hausdorff convergence}
\label{sec:conv}

\begin{theorem}[Hausdorff limit]
\label{thm:hausdorff}
Let $\Gamma = (f_0; (b_i, P_i, \varphi_i, r_i, k_i)_{i \ge 1})$ be a
countably infinite carving grammar on a compact seed $S_0$. Define
$S_i$ as above. Then
\[
  S_\infty \;:=\; \bigcap_{i=0}^{\infty} S_i
\]
is a nonempty-or-empty $G$-invariant compact set, and
$S_i \to S_\infty$ in the Hausdorff metric $d_H$ (with the convention
$d_H(\emptyset, \emptyset) = 0$ and $d_H(A, \emptyset) = \mathrm{diam}(A)$
if $A \ne \emptyset$).
\end{theorem}

\begin{proof}
By Theorem~\ref{thm:monotone}, $(S_i)$ is a nested sequence of
$G$-invariant compact sets. The intersection $S_\infty$ of a nested
sequence of compact sets in $\mathbb{R}^3$ is compact, and is
$G$-invariant because each $S_i$ is. Nested compact sequences converge
to their intersection in the Hausdorff metric (standard; see e.g.
Federer, Geometric Measure Theory, 2.10.21), which gives $S_i \to
S_\infty$.
\end{proof}

\begin{corollary}[Finite case]
For a finite carving grammar of length $n$, the result $S_n$ itself is
the Hausdorff limit of the trivial tail $(S_n, S_n, \dots)$, and is
therefore a well-defined $G$-invariant compact set.
\end{corollary}

\begin{remark}[Connection to iterated function systems]
If one restricts the Boolean op to a single carving operator applied
repeatedly with shrinking scale, the iterated carving defines a
non-autonomous contractive system on the space of compact $G$-invariant
subsets of $S_0$, and $S_\infty$ is a fixed attractor analogous to the
Sierpi\'nski fractals of classical IFS theory. The companion codebase
was seeded by exactly this observation (``SierpSphere''). The general
carving grammar generalises this to heterogeneous iteration where the
carving operator and its parameters may change at every step, and
Theorem~\ref{thm:hausdorff} is the appropriate existence statement.
\end{remark}

\section{Discussion and open questions}
\label{sec:open}

The four theorems above are structural rather than deep: they say that
a $G$-carving grammar is a well-defined monotone descent in
$\mathcal{L}_G$, that resonant grammars are a finite enumerable subset,
and that infinite carving has a limit. We list a few directions in
which one might hope to find sharper statements.

\begin{enumerate}
\item \textbf{Geometric equivalence.} Theorem~\ref{thm:finite} counts
  combinatorial types; the natural refinement is to count geometric
  equivalence classes up to $G$-action and scaling. Two distinct words
  may produce identical (or diffeomorphic, or homotopy equivalent)
  solids. A sharper classification would use the subgroup of
  $\mathcal{L}_G$ generated by resonant carves modulo the equivalence
  $\{f, f'\} \sim \{S_f = S_{f'}\}$. The image in this quotient is
  presumably still finite but its exact cardinality is open.

\item \textbf{Topological invariants along the descent.} The Euler
  characteristic $\chi(S_i)$ cannot increase arbitrarily under carving:
  removing a symmetric orbit of genus-$0$ child balls from a
  genus-$0$ seed produces a region whose $\chi$ drops by $|\Orb_G
  (\varphi_i)|$ if the orbit is cleanly separated. The question of
  whether $\chi(S_i)$ is strictly monotone, or of how the Betti
  numbers evolve, seems tractable given the explicit structure of
  carving grammars.

\item \textbf{Fractal dimension.} For a self-similar carving grammar
  --- one in which the same step is applied at a sequence of shrinking
  scales --- the Hausdorff dimension of $S_\infty$ should be computable
  via a Moran-type equation. The fitness function in the companion
  evolver already measures a discrete approximation of
  box-counting dimension, and its target value is empirically close
  to $2.3$--$2.5$. A clean theoretical statement relating the
  resonant word (in the finite alphabet of Theorem~\ref{thm:finite})
  to the fractal dimension of the limit would be very natural.

\item \textbf{Group cohomology perspective.} The assignment
  $G \mapsto \mathcal{L}_G$ is functorial in $G$. It may be profitable
  to ask what $H^*(G; \mathcal{L}_G)$ computes; perhaps the cohomology
  detects carving words that are equivalent modulo $G$-coboundaries.
  This is speculative.

\item \textbf{Dynamical systems view.} The carving grammar is a
  (heterogeneous) discrete dynamical system on $\mathcal{L}_G$ with
  attractor $S_\infty$. Standard questions: is the map expanding in
  the Hausdorff metric? Does it have a natural invariant measure?
  Is there a symbolic dynamics on the finite alphabet of
  Theorem~\ref{thm:finite} that conjugates to the geometric system?
\end{enumerate}

None of these questions are answered here. The contribution of this
note is the clean formalism of Sections \ref{sec:grammars}--\ref{sec:conv},
which makes them askable.

\end{document}
